\section{Tensor powers in the original test-configuration section module}

The polarization in a test configuration includes an isomorphism of its
restriction away from the central fibre with the pullback of the original
line bundle.  The following constructions use that degree-one isomorphism;
no comparison in higher degree, of section norms, or of metrics is assumed.
They supply the tensor-power step in the section argument behind
\cite[Lemma~5.22]{BJ22} and
\source{fhx-envelope-20261028:prop:ample-envelopes,bj22-trivval:lem:envline}.

\begin{lemma}[Canonical tensor-power transport]
\label{lem:tc-seven-canonical-power-transport}
\source{fhx-envelope-20261028:prop:ample-envelopes,bj22-trivval:lem:envline}
Let $f:Y\to X$ be a scheme morphism and $M$ a sheaf of
$\mathcal O_X$-modules.  For every $m\geq0$ there is a canonical
isomorphism
\[
 f^*(M^{\otimes m})\simeq(f^*M)^{\otimes m}.
\]
For any isomorphism $e:A\simeq B$ of module sheaves on one scheme,
its tensor powers $e^{\otimes m}$ induce linear bijections on
sections over every open.  These bijections commute with restriction
and carry $s\otimes t$ in degree $m+1$ to
$e^{\otimes m}(s)\otimes e(t)$.
\lean{EnvelopeProperty.TestConfigurationSectionTransport.sectionsEquiv,EnvelopeProperty.TestConfigurationSectionTransport.tensorPowerPullbackIso,EnvelopeProperty.TestConfigurationSectionTransport.tensorPowerIso,EnvelopeProperty.TestConfigurationSectionTransport.bundleTensorIso_tensorSection,EnvelopeProperty.TestConfigurationSectionTransport.powerSectionsEquiv,EnvelopeProperty.TestConfigurationSectionTransport.powerSectionsEquiv_restrict,EnvelopeProperty.TestConfigurationSectionTransport.powerSectionsEquiv_mul}
\leanok
\end{lemma}
\begin{proof}
The zeroth tensor power is the structure sheaf, so the canonical unit
isomorphism for pullback gives the assertion in degree zero.  In degree
$m+1$, first use the canonical isomorphism
$f^*(M^{\otimes m}\otimes M)\simeq
f^*(M^{\otimes m})\otimes f^*M$, then tensor the inductively
constructed isomorphism with the identity of $f^*M$.
All tensor products here are the sheafifications of the actual
objectwise tensor presheaves.

For $e$, start with the identity of the structure sheaf in degree zero
and recursively tensor $e^{\otimes m}$ with $e$.  Evaluation of the
resulting sheaf isomorphism and its inverse on an open gives mutually
inverse linear maps.  Naturality of the sheaf morphism proves
compatibility with restriction.  To check multiplication, represent
the section product by its actual pure tensor in the tensor
presheaf.  Naturality of the sheafification unit sends this pure
tensor to the pure tensor of the two images, which proves the stated
identity.
\end{proof}

\begin{corollary}[All original polarization powers on the original open]
\label{cor:tc-seven-original-power-identification}
\source{fhx-envelope-20261028:prop:ample-envelopes,bj22-trivval:lem:envline}
Let $j:U\to\mathcal X$ and $p:U\to X$ be scheme morphisms and let
$e:j^*\mathcal L\simeq p^*L$ be the original degree-one
polarization isomorphism.  For every $m\geq0$ this data canonically
identifies
\[
 j^*(\mathcal L^{\otimes m})\simeq p^*(L^{\otimes m}),
\]
and hence identifies their whole section modules over every open of
$U$.  Here $\mathcal L^{\otimes m}$ and $L^{\otimes m}$ are the
original tensor powers, including degree zero.
\uses{lem:tc-seven-canonical-power-transport}
\lean{EnvelopeProperty.TestConfigurationSectionTransport.polarizationPowerIso,EnvelopeProperty.TestConfigurationSectionTransport.polarizationPowerSectionsEquiv}
\leanok
\end{corollary}
\begin{proof}
Compose the canonical pullback comparison for $j$ with
$e^{\otimes m}$ and the inverse canonical pullback comparison for
$p$.  Evaluation of this isomorphism on an arbitrary open gives the
asserted linear bijection.  Each of the three factors is constructed
from the indicated morphisms and the single original isomorphism $e$.
\end{proof}

\begin{lemma}[Clearing the original parameter in every degree]
\label{lem:tc-seven-power-section-denominators}
\source{fhx-envelope-20261028:prop:ample-envelopes,bj22-trivval:lem:envline}
Let $\pi:\mathcal X\to\mathbf A^1_k$ be the original projective
test-configuration morphism, let $\mathcal L$ be an invertible module
sheaf, and write $U=\pi^{-1}D(t)$.  For every $m\geq0$, actual
restriction of the whole section module of $\mathcal L^{\otimes m}$
to $U$ is localization at the same original base parameter $t$.
In particular, for every $s\in H^0(U,\mathcal L^{\otimes m})$
there are $n\geq0$ and
$\widetilde s\in H^0(\mathcal X,\mathcal L^{\otimes m})$ such that
$t^n s=\widetilde s|_U$.
\uses{thm:scheme-module-whole-principal-localization}
\lean{EnvelopeProperty.TestConfigurationSectionTransport.TestConfiguration.parameterSection,EnvelopeProperty.TestConfigurationSectionTransport.TestConfiguration.powerRestriction_isLocalizedModule,EnvelopeProperty.TestConfigurationSectionTransport.TestConfiguration.exists_power_section_extension,EnvelopeProperty.TestConfigurationSectionTransport.TestConfiguration.exists_unitOpen_power_section_extension}
\leanok
\end{lemma}
\begin{proof}
Projectivity over the affine line makes $\mathcal X$ quasi-compact
and separated.  Every tensor power of an invertible sheaf is
invertible and therefore quasi-coherent.  The whole principal-open
localization theorem applies to this same sheaf and same parameter.
Its surjectivity condition says that every section is the
restriction of an original global section after multiplication by
an element of the submonoid generated by $t$.  Each such element is
$t^n$ for some natural number $n$, which proves the final assertion.
\end{proof}

\begin{lemma}[True model norms of the original generic tensor powers]
\label{lem:tc-seven-generic-power-norm}
\source{fhx-envelope-20261028:prop:ample-envelopes,bj22-trivval:lem:envline}
Let $\mathcal X$ be proper over a ring with a bounded valuation and
let $y$ be an original relative scheme-valued point.  Let $M$ be an
invertible module sheaf and let $s$ be any section of its original
residue-field fibre.  Define $s^{[0]}$ to be the canonical unit
section of the fibre of $\mathcal O_{\mathcal X}$, and define
$s^{[m+1]}$ by the actual native tensorator from
$s^{[m]}\otimes s$.  Then $s^{[m]}$ is a section of the actual
original fibre of $M^{\otimes m}$ and
\[
 \lVert s^{[m]}\rVert_{M^{\otimes m},y}
   =\lVert s\rVert_{M,y}^{m}
 \quad(m\geq0).
\]
The formula includes the zero section and degree zero, with $0^0=1$.
\uses{thm:actual-model-line-tensor-norm}
\lean{EnvelopeProperty.TestConfigurationSectionTransport.ModelNorm.genericPowerSection,EnvelopeProperty.TestConfigurationSectionTransport.ModelNorm.modelLineNorm_genericPowerSection}
\leanok
\end{lemma}
\begin{proof}
The canonical unit comparison takes $s^{[0]}$ to $1$ in the
original residue field.  Its true model norm is therefore the
valuation of $1$, namely $1$, which equals the zeroth power on the
right.  In degree $m+1$, multiplicativity of the true model norm for
the actual native tensor section gives
\[
 \lVert s^{[m+1]}\rVert
  =\lVert s^{[m]}\rVert\,\lVert s\rVert
  =\lVert s\rVert^m\,\lVert s\rVert
  =\lVert s\rVert^{m+1}.
\]
Induction proves the assertion.  The argument uses no division or
logarithm and therefore also applies when $s=0$.
\end{proof}

\begin{remark}
These statements identify genuine tensor powers and clear their original
section denominators.  To deduce the literal formula in
\source{fhx-envelope-20261028:prop:ample-envelopes}, one still needs
the Laurent expansion of sections on $X\times\mathbf G_m$ with
coefficients in $H^0(X,L^{\otimes m})$, the comparison of its actual
Gauss norm with the original model metric, the common-model
competitor comparison, eventual verticality, and the pointwise Fekete
limit.  None of these is asserted by the preceding lemmas.
\end{remark}
